\section{PD-TDM Design}
\label{sec:design}

PD-TDM runs inside one model instance on the same TP devices. Weights and KV
blocks remain local. Before every model iteration, the scheduler selects phase
$P$ or $D$. A $P$ iteration advances only unfinished prompts; a $D$ iteration
generates one token for each selected active sequence. This phase-purity
invariant is enforced for running partial prefills and newly admitted requests.

\textbf{Bounded continuation.}
The prefill planner applies an aggregate $C=2{,}048$ token limit per $P$
iteration. It first continues partial prefills and then admits waiting requests
with the remaining budget. An unfinished prompt retains its KV state and
continues later. The token bound limits admitted work; since iteration time is
workload dependent, it is not a strict wall-clock guarantee.

\textbf{Phase selection.}
A token bucket accumulates a configurable target prefill share $r\in[0,1]$.
When both phases have work, sufficient credit selects $P$ and consumes one unit;
otherwise it selects $D$. If only one phase has work, the scheduler immediately
serves it, so $r$ describes a contention-time tendency rather than a fixed
sequence or wall-clock ratio. The implementation also exposes TTFT-urgency and
decode-starvation overrides. Our claims concern the observed configured
scheduler, not an online controller that finds an optimal $r$.

This design predicts its boundary. When prompt queueing dominates and TPOT has
slack, concentrated $P$ iterations can move TTFT left while $D$ iterations
remain within the TPOT threshold. Under a deep decode queue, the silence created
by a $P$ iteration consumes the remaining TPOT slack and mixed chunked prefill
can be preferable.
